\section{Uniform control of discrete endpoints}
\label{spike:section}

The preceding section supplies cheap surrounding circuits.  We now use
them to control distances from every prescribed discrete vertex.  A
circuit's length does not bound the cost of reaching it from the vertex
it surrounds.  Our deterministic covering lemma keeps these access costs
in a bounded number of shorter chronological intervals.  Rank credit
allows the covering to be iterated, and the remaining intervals are
finally traversed by the contour walk itself.

The first output is a uniform modulus for endpoints close in contour
time.  The second joins arbitrary representatives of the same limiting
spatial point, using the finite horizontal-chain relation.  Both
conclusions retain a prescribed open confinement.

Throughout this section, $v_n(j)$ is the primal tour vertex at integer time
$j$.  The continuous interpolation of $j\mapsto v_n(j)$ follows a primal
tree edge or stays at a vertex at each step.  In particular,
\begin{equation}\label{spike:chronological-bound}
 d_{G_n}\bigl(v_n(i),v_n(j)\bigr)\le |i-j|.
\end{equation}
Whenever confinement is specified, the right side is still an upper bound
if the intervening chronological walk stays in that confinement.

\subsection{Covering an interval by circuit interiors}

We use a simple circuit as an embedded Jordan curve consisting of primal
edges, with a specified disk on its inside.  In a planar chart this is the
bounded disk.  Length means the number of its edges, with multiplicity one
on the circuit.  Loops and two-edge circuits in a multigraph are allowed
when their embeddings are Jordan curves.

\begin{definition}[An available circuit]\label{spike:available}
Let $v:\mathbb Z\to V(G)$ be a nearest-neighbor walk with allowed stays in
an embedded planar graph, and interpolate it along graph edges.  Fix
integers $d\ge1$ and $\ell\ge 4d$, a positive number $L$, and an open
planar confinement $U$.  An integer time $t$ is good if there is a simple
circuit $C_t\subset U$, of length at most $L$, with interior $D_t$, such
that
\begin{enumerate}[label=(\roman*)]
 \item the walk on $[t-d,t+d]$ is contained in $D_t$;
 \item no chronological interval of $\ell$ consecutive steps is
       contained in $\overline D_t$.
\end{enumerate}
All other times are bad.  Condition (ii) is required only in the
chronological window in which the exit searches below are made.
\end{definition}

For vertices $x,y\in U$, let $d_{G,U}(x,y)$ denote the shortest-path
distance restricted to graph paths in $U$, with value infinity when no
such path exists.

\begin{lemma}[Covering by circuit interiors]\label{spike:covering}
Fix an integer $h\ge1$.  There is a constant $K_h<\infty$ with the
following property.  Let $I=[i,j]\cap\mathbb Z$ have length at most $T$.
Suppose that no $h$ bad times in $I$ are pairwise at distance at least
$\ell$, and that the interpolated walk on the $K_h\ell$ enlargement
of $I$ lies in $U$.  Then there are at most $h+2$ pairs of integer times
$(i_r,j_r)$, all in the $K_h\ell$ enlargement of $I$, satisfying
\[
 |i_r-j_r|\le K_h\ell,
\]
and graph paths in $U$ of total length at most
\[
 K_h(1+T/d)L,
\]
which, together with connections between $v(i_r)$ and $v(j_r)$, form a
connection between $v(i)$ and $v(j)$.  Consequently,
\begin{equation}\label{spike:covering-recursion}
 d_{G,U}(v(i),v(j))\le K_h(1+T/d)L
       +\sum_{r=1}^{h+2}d_{G,U}(v(i_r),v(j_r)).
\end{equation}
Unused summands can be taken to be zero.  The conclusion also holds when
there are no good times in $I$.
\end{lemma}

\begin{proof}
Choose a maximal $\ell$-separated set of bad times in $I$.  It has at
most $h-1$ elements.  The union of their chronological $\ell$-neighborhoods
covers every bad time.  Merge overlapping neighborhoods and intersect
with $I$.  Their total length is at most $2(h-1)\ell+2h$, and the
remaining set has at most $h$ integer intervals, all of whose times are
good.

Consider one such interval $[a,b]\cap\mathbb Z$.  Choose good times
$a=t_1<\cdots<t_k=b$ with successive differences at most $d$, using at
most $2+(b-a)/d$ times.  For each one choose an available circuit.  The
interiors of consecutive circuits overlap: the walk at each time between
$t_u$ and $t_{u+1}$ belongs to both interiors on a nonempty chronological
subinterval.

Among this finite family retain the disks maximal for inclusion.  Every
discarded disk is contained in a retained disk.  The intersection graph
of retained interiors is connected: a chain of overlapping original
disks maps to a chain of overlapping containing maximal disks.  Two
distinct maximal Jordan disks which overlap have intersecting
boundaries.  Indeed disjoint Jordan boundaries with intersecting
interiors would force one disk to contain the other.  Thus the union of
the retained circuit boundaries is a connected subgraph.  Its number of
edges is at most the sum of the lengths of the original circuits.

Select a retained disk containing the first disk, so that it contains
$v(a)$ in its interior.  Condition (ii) implies that the interpolated
walk first reaches its boundary at an integer time in $[a,a+\ell+1]$.
The intersection is a graph vertex, since both the walk and the circuit
use edges of the same planar graph.  Similarly a retained disk containing
the last disk has a boundary vertex represented at a time in
$[b-\ell-1,b]$.  The connected union of retained boundaries joins these
two vertices at a cost no greater than the sum of the chosen circuit
lengths.  We have made no estimate of the cost of reaching either
boundary from $v(a)$ or $v(b)$.

Do this for every remaining good interval, and concatenate the resulting
boundary paths in the order of those good intervals.  Within a short good
interval the forward hit may follow the backward hit; orient its boundary
path from the forward hit to the backward hit regardless of their time
order.  Each gap between these paths is
represented by a pair of times separated by a removed bad interval and
at most $2\ell+2$ extra steps.  Initial and final gaps obey the same bound.
There are at most $h+1$ gaps; allowing $h+2$ covers endpoint and rounding
conventions.  Every gap has length at most $K_h\ell$, and every time used
lies in the $K_h\ell$ enlargement of $I$, for a constant depending only
on $h$.  The total number of circuits is at most
$2h+T/d$, which gives the stated cost after increasing $K_h$.

If no good interval remains, the length of $I$ is at most the total
length of the bad neighborhoods.  Use the single unresolved pair $(i,j)$.
This proves all claims, including Equation~\eqref{spike:covering-recursion}.
\end{proof}

\subsection{Rank credit bounds the separated bad times}

We next show that the obstruction to the covering lemma---many separated
times without an available circuit---cannot persist under the selected
laws used in the recursion.  All scale constants are fixed here, before
choosing such a law.

Let $s>0$ be a polynomial-ratio exponent for the crossing normalization,
as in Equation~\eqref{comp:polynomial-normalization}:
\begin{equation}\label{spike:normalization-input}
 \frac{B(u)}{B(v)}\le K(u/v)^s,\qquad 1\le u\le v.
\end{equation}
It follows, by fixing $u$, that $B(v)\ge c v^s$.  Decreasing $s$ preserves
Equation~\eqref{spike:normalization-input}; we henceforth take $0<s<1$.
Choose constants in the order
\begin{equation}\label{spike:parameter-order}
 0<b-1<s/8,\qquad s'=s/2,\qquad
 1<a'_0<a_0<a<b,
\end{equation}
where $a-1$ is sufficiently small relative to
$(d_2-1/2)(b-a)$ for the rank-credit argument.  The integer $h$ is then
chosen sufficiently large for that argument.  All these quantities
depend on the fixed $q$, but not on $n$.

Here $b>1$ is a log-window constant and is unrelated to the inventory
parameter $b_F$.  At level $N$, put, with integer rounding understood,
\begin{equation}\label{spike:scales}
 T=n e^{-N},\qquad \ell=n e^{-a_0N},\qquad
 d=n e^{-bN},\qquad L=B(n)e^{-s'N}.
\end{equation}

We use the preselection events of Lemma~\ref{view:macroscopic-cutoffs}.
Fix a compact chronological window, a spatial compactness and clearance
cutoff, an error tolerance $\epsilon>0$, and a deterministic largest
view radius $\rho_n n$ with $\rho_n\to0$.  Before selecting a bad level
or its bases, that lemma constructs a master event $\mathcal T_n$ and
a deterministic depth bound $D_n\to\infty$.  Its failure probability
on the fixed compactness and clearance cutoff has limsup at most
$2\epsilon$.  On $\mathcal T_n$, simultaneously for every depth
$d\le D_n$, all explorations from the tested window with at most $d$
steps, each a scalar link or a jump of size at most $d\rho_n n$, lie in the
corresponding
deterministic chronological enclosure and have the required spatial
confinement.  The event uses macroscopic contour barriers and uniform
finite-chain projection control; it is chosen under the unselected law.

After conditioning on a bad configuration, the finite-list depth,
enclosure, and coding cutoffs may be increased only along a slower
diagonal within these preselected bounds.  In particular, the sparse
coding conditions of Lemma~\ref{view:all-seed-code} are imposed by
slowing that diagonal.  An ordinary high-probability enclosure is never
transferred through a subsequently chosen density $e^{o(N)}$.

\begin{proposition}[Circuit sparsity input]\label{spike:circuit-sparsity}
Use the constants above and the available-circuit definition with the
scales in Equation~\eqref{spike:scales}.  Let $E_{n,N}$ be the event that
some chronological interval of length at most $T$, in the tested compact
window, contains $h$ pairwise $\ell$-separated bad times.  Along every
sequence for which
\[
 N\longrightarrow\infty,\qquad n e^{-bN}\longrightarrow\infty,
\]
there is no sequence of probability laws $\mathbb Q_{n,N}$ supported on
$E_{n,N}$ and the preselected master events just described, with
$e^{-N}\le\rho_n\to0$, such that
\begin{equation}\label{spike:selected-law}
 \mathbb Q_{n,N}\le \exp(\kappa_{n,N})\mathbb P_n,
 \qquad \kappa_{n,N}=o(N).
\end{equation}
The same assertion holds for any fixed constant multiple of $T$ and
for sufficiently small fixed spatial neighborhoods inside the selected
confinement.
\end{proposition}

\begin{proof}
Mark the $h$ selected bad bases and finite witness
diagrams before choosing the independent shifted address grids.
The address estimate of Section~\ref{arm:section} bounds the integrated
box count under Equation~\eqref{spike:selected-law} and charges
$\beta$ times the integrated monochrome rank, up to arbitrarily small
relative errors.  Proposition~\ref{view:acyclic-circuit} shows that a
positive probability of an acyclic base component at a uniform viewing
level in $[aN,bN]$ would supply an available circuit at one of the bad
bases.

We relate these center components to the projected graphs charged by
the address estimate.  The bases are separated by
$\ell/r(X)\ge e^{(a-a_0)N}$ at a view level $X\in[aN,bN]$, so their
enrolled centers are distinct.  On each fixed finite truncation,
Lemma~\ref{view:returns}(ii) supplies actual links witnessing a
spanning tree of every scalar center class.
Proposition~\ref{view:saturation} identifies these links with the
graph attained by the seed recipe.  The all-seed code includes its
actual witnesses and the jumps associating their endpoints to their
frozen centers.  Outside the levels charged by the integrated jump
error, these endpoints and centers occupy the same projected cells;
distinct center lines occupy distinct cells.  Bases also have tight
seed offsets from their assigned centers.  Apply the independent-grid
boundary estimate to the original centers, chosen before that grid:
conditional on a bounded seed offset, the crossing probability is at
most a fixed multiple of $\min\{1,e^{-U_j}\}$, whose uniform-coarsening
average tends to zero.  Rounding an origin keeps it in the same
terminal cell, and hence in the same cell at every coarser generation.
Thus each base and its frozen origin occupy the same projected cell
with probability tending to one.

Thus a rank cycle of a limiting base component is witnessed inside
the corresponding projected base component.  Adding connected vertices
cannot destroy it: extend monochrome spanning forests by a union
spanning tree rooted at the old component, gaining at least one unit
of total rank per new vertex.  All other code entries may therefore
be kept.  The acyclic-circuit estimate makes the probability of an
acyclic base component at a uniform view level tend to zero.  Averaging
over the seed-generation trials and applying Markov's inequality now
shows that every projected base component has a rank cycle outside
$o(N)$ generations, in probability.  Fixed truncations are removed
along the sparse diagonal of Lemma~\ref{view:all-seed-code}.

The separation hypothesis of Proposition~\ref{arm:rank-credit}
holds for every displacement.  At a fixed chronological cutoff, the
largest charged cell width is $O(ne^{-aN})=o(\ell)$.  Along the
permitted enlarging cutoff, the all-seed condition
$\Lambda_j\log(2\widehat C_j)/L_j\to0$, together with $L_j=o(N_j)$,
gives $\log\widehat C_j=o(N_j)$, so the same comparison remains true.

The rank-credit estimate gives negative credit at least
$(d_2-1/2-o(1))(b-a)hN$ on this range, whereas the earlier possible
positive credit is at most $C_0N+C_1(a-1)hN+o(N)$.  The choices in
Equation~\eqref{spike:parameter-order}, followed by the choice of $h$,
make the former strictly greater.  The selected density costs only
$o(N)$, so this contradicts the address probability bound and proves
the proposition.
\end{proof}

\begin{remark}\label{spike:selection-strength}
The bounded-density hypothesis in Equation~\eqref{spike:selected-law}
is intentional.  Conditioning on an event of probability
$\exp(-o(N))$ has precisely this property, and that is the only kind of
selection used below.  A bound on relative entropy alone would transfer
vanishing reference probabilities, but would not by itself give the
uniform first-moment bounds for arbitrarily large address systems used
in constructing the views.  We do not make that stronger inference.
\end{remark}

\subsection{Iteration with a growing lattice cutoff}

The recursion pays for circuit boundaries at each stage and passes only
the short access intervals to the next stage.  We must ensure that all
stages are valid together, while stopping before the lattice estimates
reach a bounded time scale.

For a chronological window $J$ and spatial confinement $U$, write
\[
 \omega_n(N;J,U)=\frac1{B(n)}
 \sup_{\substack{i,j\in nJ\cap\mathbb Z\\ |i-j|\le n e^{-N}}}
 d_{G_n,U}\bigl(v_n(i),v_n(j)\bigr).
\]
Window enlargements below are measured in units of $n$.  On the
complement of $E_{n,N}$, Lemma~\ref{spike:covering}, applied with
Equation~\eqref{spike:scales}, gives for sufficiently large $N$
\begin{equation}\label{spike:omega-recursion}
 \omega_n(N;J,U)
 \le C_h e^{-cN}+H\,
 \omega_n(a'_0N;J^{C_h e^{-a_0N}},U),
 \qquad c=s'-(b-1)>0,
\end{equation}
where $H=h+2$.  Indeed
$K_h\ell\le n e^{-a'_0N}$ for sufficiently large $N$, and the circuit
payment is at most
$K_h(1+e^{(b-1)N})e^{-s'N}$.  Integer rounding is absorbed in the same
constants because the smallest scale $d$ tends to infinity.

\begin{lemma}[Uniform validity along the recursion]\label{spike:levels}
Assume Proposition~\ref{spike:circuit-sparsity}.  Start at a deterministic
$N_0=N_0(n)\to\infty$ and set $N_j=(a'_0)^jN_0$.  Apply the recursion
only at those levels satisfying
\[
 N_j\le(1-s/2)\log n.
\]
After the compactness and clearance cutoffs have been exhausted, the
probability that any applied level has a bad configuration tends to zero.
\end{lemma}

\begin{proof}
Suppose the probability of some bad level is bounded below by
$\delta>0$ along a subsequence on a fixed compactness and clearance
cutoff.  Choose $\epsilon<\delta/8$ and apply
Lemma~\ref{view:macroscopic-cutoffs} with $\rho_n=e^{-N_0}$ to a fixed
slight enlargement of the tested window.  This gives one master event
$\mathcal T_n$ for every applied level, before any level or bases are
selected.  The union of the bad-level events, intersected with
$\mathcal T_n$ and the fixed cutoff, still has probability at least
$\delta/2$ for all large $n$.
Since $\sum_{j\ge0}(j+1)^{-2}<\infty$, there is a deterministic
index $j=j(n)$ for which the intersection of that level's bad event
with these same preselected events has probability at least
$c\delta(j+1)^{-2}$.  Condition on this intersection and then choose the
bad bases by the finite lattice ordering.  The logarithm of the
density bound is at most $C_\delta+2\log(j+1)$, whereas
$N_j=(a'_0)^jN_0$; thus the density bound is $\exp(o(N_j))$.

At every applied level,
\[
 n e^{-bN_j}\ge n^{1-b(1-s/2)}\longrightarrow\infty,
\]
since $0<s<1$ and $b-1<s/8$ imply
\[
 b(1-s/2)<(1+s/8)(1-s/2)<1.
\]
All interval-window enlargements fit in a fixed slightly larger window:
their total size is bounded by
$\sum_j C_h e^{-a_0(a'_0)^jN_0}$, which tends to zero.
The resulting conditional laws contradict
Proposition~\ref{spike:circuit-sparsity}.  Letting the arbitrary
preselection loss tend to zero and exhausting the compactness and
clearance cutoffs proves the claim.
\end{proof}

\begin{proposition}[Chronological endpoint modulus]\label{spike:time-modulus}
Assume Equation~\eqref{spike:normalization-input} and
Proposition~\ref{spike:circuit-sparsity}.  In a bilateral contour coupling,
let $J$ be a compact time interval whose limiting curve image has a
neighborhood compactly contained in an open set $U$.  Then, for every
deterministic $\delta_n\downarrow0$,
\[
 \frac1{B(n)}
 \sup_{\substack{i,j\in nJ\cap\mathbb Z\\ |i-j|\le n\delta_n}}
 d_{G_n,U}\bigl(v_n(i),v_n(j)\bigr)
 \longrightarrow0
 \quad\hbox{in probability}.
\]
The assertion also holds locally on the event that the stated
compact-containment condition is satisfied, by exhausting its clearance
and the chronological and spatial compactness cutoffs.
\end{proposition}

\begin{proof}
The claim is immediate when $n\delta_n<1$, since only identical integer
times are tested.  Otherwise set $N_0=-\log\delta_n$.
If $N_0>(1-s/2)\log n$, chronological steps
give a bound $C n^{-s/2}$ after normalization.  Otherwise iterate
Equation~\eqref{spike:omega-recursion} until the first new level
$N_{j_*}>(1-s/2)\log n$.  Lemma~\ref{spike:levels} shows that all recursions
are valid with probability tending to one.  Their total normalized
payment is bounded by
\[
 \sum_{j\ge0}H^j C_h e^{-c(a'_0)^jN_0},
\]
which tends to zero.  For example, after increasing $N_0$, the ratio of
successive terms is at most $1/2$, so the sum is at most twice its first
term.

Each terminal unresolved pair has chronological separation at most
$n e^{-N_{j_*}}\le n^{s/2}$.  There are at most $H^{j_*}$ such pairs, and
$j_*\le C_1+C_2\log\log n$.  Equations
\eqref{spike:chronological-bound} and
\eqref{spike:normalization-input} bound their total normalized cost by
\[
 C(\log n)^{C_3}n^{-s/2}\longrightarrow0.
\]
Their chronological paths stay in $U$ with probability tending to one
by the uniform contour coupling and the stated clearance.  The same is
true of all earlier circuit paths by their definition.  This proves the
assertion, including confinement.
\end{proof}

\subsection{Joining endpoints at the same limiting point}

Close spatial endpoints need not have close contour times.  The sewing
relation connects their limiting times by finitely many horizontal
generators.  Approximate each generator by a genuine discrete link, then
use the chronological modulus to join its small endpoint errors.

\begin{proposition}[Local endpoint joining]\label{spike:spatial-modulus}
Assume Proposition~\ref{spike:time-modulus},
Lemma~\ref{bg:chord-approximation}, and the compact sewing interface
(T1)--(T3) of Section~\ref{bg:topology-section}, with vanishing projected
edge diameters.  In a joint extraction on such a chart, arbitrary
primal-vertex endpoints whose continuum locations approach the same point
can be joined at cost $o(B(n))$, with the path confined to any prescribed
open neighborhood of that point.  The conclusion is uniform over
endpoints in a fixed compact chart, in the usual sequential sense.
The same conclusion applies in a bilateral exhaustion whenever the
tested vertices have tight chronological representatives in its
compact sewings.
\end{proposition}

\begin{proof}
If uniformity failed, select a violating pair with positive probability
and then a further coupled subsequence.  First restrict to a compact
sewing; in an exhaustion, truncate the stipulated tight chronological
representatives.  This truncation is part of the hypothesis; it is not
deduced from a passage comparison.  Property (T1) gives limiting representative times
$t$ and $u$ with $\eta(t)=\eta(u)$.  Their equivalence is generated by a
finite chain of horizontal coordinate identifications by (T3).  For each
ordinary generator, strict-height approximations yield actual primal
identifications or genuine primal edges whose time representatives
approach its two endpoints.  At a branch visit one first splits into
the finitely many consecutive ordinary generators.  This is a
topological statement about the discrete sewing; it estimates no
chronological connector cost.

Choose closed rational time neighborhoods of all the finitely many
visits whose curve images lie in the prescribed open spatial
neighborhood.  Truncate the number of visits and then their countable
rational neighborhood choices to a finite list, with arbitrarily small
probability loss.  The finitely many time errors in the link approximations
are then bounded, with probability tending to one, by a deterministic
sequence $\delta_n\downarrow0$, chosen sufficiently slowly.
Proposition~\ref{spike:time-modulus} joins the selected original
representatives to the approximating link endpoints, and joins
successive approximating endpoints at each visit, at cost $o(B(n))$.
The finitely many genuine edge costs are also $o(B(n))$ because
$B(n)\to\infty$, and their projected diameters tend to zero, so they
stay inside the prescribed neighborhood.  Concatenation contradicts
the selected violation after the arbitrary truncation losses vanish.
This proves the uniform sequential assertion.  Notice that a
second-coordinate scalar identification is not itself declared a
unit-cost primal edge: it is the strict-height matched-step
approximation, together with the chronological modulus just proved,
that supplies the primal connection.
\end{proof}

\begin{corollary}[Bilateral endpoint modulus]
\label{spike:endpoint-modulus}
Under their respective hypotheses, the chronological and spatial
endpoint conclusions of Propositions~\ref{spike:time-modulus} and
\ref{spike:spatial-modulus} hold simultaneously.
\end{corollary}

The universal endpoint tests recorded in the certificate framework
allow this conclusion to enter the common-kernel transfer.  It is not
inferred by reversing a semicontinuity assertion for existential path
certificates.
